\section{Literature Review}
\label{sec:lit}
Our work intersects the literature on realized volatility forecasting, the econometric evaluation of predictive accuracy, and the application of deep learning to tabular financial data. 

The baseline for our analysis is the HAR-RV model introduced by \cite{corsi2009simple}, which elegantly models the volatility cascade generated by heterogeneous market participants. While standard loss functions like Mean Squared Error (MSE) are common, \cite{patton2011volatility} demonstrates that the QLIKE loss function is robust to noise in the volatility proxy, making it the preferred metric for ranking volatility forecasts. To rigorously assess forecast superiority, we rely on the Diebold-Mariano test \cite{diebold1995comparing} with the small-sample correction of \cite{harvey1997testing}, as well as the Model Confidence Set (MCS) framework of \cite{hansen2011model}, which isolates a set of models that contain the best model at a given confidence level.

In recent years, the literature has increasingly explored the integration of sentiment into volatility models. The advent of domain-specific language models, such as FinBERT \cite{araci2019finbert, yang2020finbert}, has provided researchers with robust, continuous sentiment indices. However, feeding high-dimensional, noisy sentiment data directly into linear autoregressive models often yields poor out-of-sample predictions. To address this, we draw on recent advances in deep learning for tabular data, specifically the Feature Tokenizer Transformer (FT-Transformer) and Gated Residual Networks \cite{gorishniy2021revisiting}. By combining a linear prior with a deep learning residual branch, our Neural-HAR model seeks to overcome the linear constraints of classical models while avoiding the optimization instability of pure neural networks in low signal-to-noise environments.
